\section{Product scope and user experience}

\subsection{Connected application surfaces}
\begin{table}[htbp]
\centering
\small
\begin{tabularx}{\textwidth}{@{}p{31mm}X X@{}}
\toprule
\textbf{Surface} & \textbf{Purpose} & \textbf{Important behavior} \\
\midrule
Explore & Choose a launch city & Bundled packs; populated, empty, loading, and
error states; no network dependency \\
Destination & Inspect a city and choose mode & Curated route cards, Free Run,
pack version, landmarks, and quests \\
Route Details & Understand one route & Distance, duration, difficulty, safety
notes, quest list, and explicit start \\
Run Tracker & Record the activity & Permission explanation; start, pause,
resume, finish, discard, checkpoints, and quest actions \\
Quest Camera & Capture evidence & Proximity or fallback validation, retained
photo, required caption, cancellation and camera-error recovery \\
Run Summary & Review the result & Distance, duration, pace, track, evidence,
achievements, Story Studio, and optional online sharing \\
Journey & Reopen local work & Runs, editable stories, achievements, refresh,
deletion confirmation, and empty/error states \\
Story Studio & Compose the visual diary & Original templates, transforms,
layer order, undo/redo, local save/reopen, exact PNG export, share sheet \\
\bottomrule
\end{tabularx}
\caption{Production application surfaces and their responsibilities.}
\end{table}



The production navigation provides more than the minimum three to four connected
screens. The bottom navigation keeps Explore and Journey understandable, while
route selection, tracking, evidence capture, summary, and Story Studio appear as
focused tasks in the natural order of use.

\begin{figure}[htbp]
\centering
\begin{tikzpicture}[
  node distance=5mm and 6mm,
  box/.style={draw=QuestInk,rounded corners=2mm,fill=QuestPaper,
              minimum width=29mm,minimum height=10mm,align=center,font=\small},
  arr/.style={-{Latex[length=2mm]},thick,color=QuestTeal}]
  \node[box] (explore) {Explore};
  \node[box,right=of explore] (city) {Destination};
  \node[box,right=of city] (route) {Route Details};
  \node[box,below=of city] (tracker) {Run Tracker};
  \node[box,left=of tracker] (camera) {Quest Camera};
  \node[box,right=of tracker] (summary) {Run Summary};
  \node[box,right=of summary] (studio) {Story Studio};
  \node[box,below=of summary] (journey) {Journey};
  \draw[arr] (explore) -- (city);
  \draw[arr] (city) -- (route);
  \draw[arr] (route) -- (tracker);
  \draw[arr] (city) -- node[right,font=\scriptsize]{Free Run} (tracker);
  \draw[arr] (tracker) -- (camera);
  \draw[arr] (camera) -- (tracker);
  \draw[arr] (tracker) -- (summary);
  \draw[arr] (summary) -- (studio);
  \draw[arr] (summary) -- (journey);
  \draw[arr] (journey) -- (summary);
  \draw[arr] (journey) to[bend left=24] (studio);
\end{tikzpicture}
\caption{Questory's end-to-end production user flow.}
\end{figure}



\subsection{End-to-end user flow}
\begin{table}[htbp]
\centering
\footnotesize
\begin{tabularx}{\textwidth}{@{}p{7mm}p{31mm}p{52mm}X@{}}
\toprule
\textbf{Step} & \textbf{User action} & \textbf{Application response} &
\textbf{Offline and recovery behavior} \\
\midrule
1 & Launch Questory & Initialize local dependencies and open Explore with the
Nha Trang and Ho Chi Minh City packs. & If initialization fails, show a retryable
local-storage error instead of an empty or misleading home screen. \\
2 & Choose a destination & Display bundled routes, landmarks, quests, estimates,
and safety notes. & All launch content is packaged with the application; map tiles
and a backend are not required. \\
3 & Select a curated route or Free Run & Carry the selected route into tracking,
or create an unrestricted session using the same recording pipeline. & The user
can return without starting; no data is recorded before an explicit start. \\
4 & Start tracking & Explain location use, check GPS service, request permission,
and begin checkpointed recording after approval. & Denial and disabled GPS lead
to recovery guidance; pause, resume, finish, and discard remain explicit. \\
5 & Complete a photo quest & Validate proximity or a documented fallback, open
the camera, require evidence and any necessary caption, then return to tracking.
& Captured evidence is copied from temporary camera storage into app-controlled
local storage. \\
6 & Finish the run & Derive duration, distance, pace, quest results, and personal
achievements; persist the summary; then offer Story Studio or Journey. & Invalid
GPS jumps are filtered, zero-distance pace is handled, and the completed run is
available without a network. \\
7 & Build and share a story & Create or reopen an editable 1080 by 1920 project,
arrange content, save it, export a deterministic PNG, and open Android sharing.
& Editing and export are local; a failed share does not remove the saved project
or completed run. \\
8 & Revisit Journey & Browse prior runs and stories, reopen a summary or project,
or deliberately delete a run. & Deletion removes related stories first, restores
them if run deletion fails, and cleans retained photos only after consistency is
confirmed. \\
\bottomrule
\end{tabularx}
\caption{Detailed production user flow and its recovery paths.}
\end{table}


The diagram and table describe the same production journey at two levels: the
diagram shows how screens connect, while the table records the user action,
application response, and recovery behavior at each stage. Importantly, neither
curated routes nor Free Run depend on a server. Bundled destination content,
SQLite persistence, device location, retained photographs, story editing, and
PNG export remain available offline. Only explicitly optional integrations,
such as live sharing, may require connectivity.

\subsection{Primary journey in detail}
The user begins in Explore, where both launch packs are loaded from versioned
JSON bundled with the application. Selecting a city reveals its routes,
landmarks, quests, pack version, and offline status. A curated route provides an
expected path, estimated distance, duration, and safety notes. Free Run keeps the
same recording and storytelling pipeline without requiring a route.

Immediately before tracking, Questory explains why precise location is needed.
After confirmation, the application checks the device service and requests
permission. The tracker exposes start, pause, resume, finish, and discard actions.
Quest cards display their radius and fallback rule. Capture opens the device
camera only after proximity or explicit fallback confirmation.

Finishing produces a summary from the immutable session: active duration,
distance, pace, track, landmarks, completed and skipped quests, retained evidence,
and derived achievements. The result is stored in History and may become the
source of a Story Studio project. The story can be edited, saved, reopened,
exported as a 1080 by 1920 PNG, and passed to the Android share sheet.

\appscreenshot{images/explore.png}{102mm}{Explore uses bundled city packs and clearly communicates offline availability.}

\subsection{Interface direction}
The visual language uses paper-like neutral surfaces, teal actions, coral and
yellow accents, heavy display type, rounded cards, film frames, labels, and
stickers. Solid colors avoid reproduction surprises and satisfy the product rule
against gradients. Large controls support outdoor use; status labels and text do
not rely on color alone.
